\section{能源约束与硬件迭代：性能竞争的真实边界}

\subsection{计算的物理约束：最终是能量预算}

Jensen 在中后段明确指出：计算扩展受限于能源与热。无论是训练还是 inference，本质都要在功耗、散热、延迟与吞吐之间做折中。

\begin{knowledgebox}{为什么“性能提升”越来越依赖系统设计}
单纯提升单芯片频率已经不足以支撑 AI 工作负载，必须依赖：
\begin{itemize}
    \item 芯片架构优化（并行度、内存层级）
    \item 高速互联（跨 GPU/节点通信）
    \item 系统级软件调度（编译器、runtime、并行策略）
\end{itemize}
因此当前竞争是“系统工程竞争”，而非“单器件竞争”。
\end{knowledgebox}

\subsection{从数据中心到个人设备：AI supercomputer 的形态下沉}

访谈里 Jensen 展示了早期系统，并提到从高价研究设备向更低门槛个人设备演进。这反映了同一趋势：AI 计算会像当年的个人计算一样下沉到更广泛开发者与学生群体。

\begin{figure}[H]
\centering
\includegraphics[width=0.88\textwidth]{frame-dgx-a.jpg}
\caption{访谈中回顾早期 AI supercomputer 交付场景 \protect\footnotemark}
\end{figure}
\footnotetext{视频画面时间区间：00:36:58--00:37:06。}

\begin{figure}[H]
\centering
\includegraphics[width=0.88\textwidth]{frame-geforce-a.jpg}
\caption{主持人展示可桌面化部署的个人 AI supercomputer 形态 \protect\footnotemark}
\end{figure}
\footnotetext{视频画面时间区间：00:52:33--00:52:45。}

\begin{importantbox}{关键结论}
“AI democratization”不仅是模型开放，更是计算形态和开发工具链向个人可及范围迁移。教育与产业创新速度将因此发生二次加速。
\end{importantbox}

\subsection{本章小结}
AI 竞争已从“谁有更大模型”升级为“谁能在能源约束下持续提供更高系统效率与更低使用门槛”。

